% Template LaTeX: Standar Sections
% Gunakan file ini sebagai template untuk membuat section baru

\section{Your Section Title Here}

\subsection{Subsection Title}

Tulis konten Anda di sini. Gunakan struktur berikut:

\subsubsection{Subsubsection (if needed)}

Dengan struktur modular ini, Anda dapat:
\begin{itemize}
    \item Fokus pada satu section tanpa distraksi
    \item Dengan mudah collaborate dengan tim member lain
    \item Maintain version control yang lebih baik
\end{itemize}

\subsection{Another Subsection}

\textbf{Tips:}
\begin{enumerate}
    \item Gunakan descriptive section titles
    \item Simpan images di direktori \texttt{images/}
    \item Tambahkan references ke \texttt{bib/custom.bib}
    \item Cite dengan \texttt{\textbackslash cite\{key\}}
\end{enumerate}

% Contoh Figure
\begin{figure}[h]
  \centering
  \includegraphics[width=\columnwidth]{images/example.png}
  \caption{Your caption here}
  \label{fig:example}
\end{figure}

% Contoh Table
\begin{table}[h]
  \centering
  \caption{\label{tab:example}
    Your table caption here.
  }
  \begin{tabular}{lll}
    \hline
    \textbf{Column 1} & \textbf{Column 2} & \textbf{Column 3} \\
    \hline
    data1 & data2 & data3 \\
    \hline
  \end{tabular}
\end{table}

% Note: Jangan lupa tambahkan \input{sections/your_file.tex} di main.tex!
